\chapter{Introduction}
\label{ch:intro}

\section{Scope of the work}
\label{sec:intro-scope}

This work examines the acceleration of top-k selection, that is, finding the k
largest elements of a set of n values, on three architectures: a
general-purpose processor, a graphics unit and a neural processing unit
integrated into the same die as the processor.

For each architecture, two algorithmic approaches with radically different
execution characteristics are developed. Truncated bitonic sort executes a
fixed sequence of comparisons independent of the data, while selection with
local heaps under the map-reduce scheme adapts its work to the values of the
input. All six implementations are measured over the same parameter space, with
a common harness and common measurement semantics, so that the differences that
emerge are owed to the algorithms and the hardware, not to the way of
measuring.

The central question is under what conditions offloading to an accelerator pays
off, weighing both the speed of the architecture and the additional cost of the
process.

\section{Motivation}
\label{sec:intro-motivation}

Neural processing units are now integrated on the same die as the processor of
modern laptops, share its main memory and are available to any application that
can make use of them. Whether it is worth offloading this computation to the
accelerator has become a timely question for developing software on such
hardware.

They were, however, designed for a specific workload. Their programming model
is organised mainly around neural-network operators, and their architecture
around the assumption that every piece of data will be used many times after it
is transferred. What happens when the workload does not have this shape remains
largely unexplored.

Top-k selection is a suitable problem for posing this question. It is simple to
state and widespread in practice, as it appears in nearest-neighbour search, in
retrieval and recommendation systems, in the sampling of language models and
many other applications. At the same time, accelerators were designed for
workloads with many operations per element and extensive data reuse, whereas
top-k selection has minimal operations per element and no reuse at all.

\section{Contribution}
\label{sec:intro-contribution}

\paragraph{Implementations.} Six implementations were built, two algorithms on
three architectures, with support for five data types. On the processor, the
AVX-512 vector extensions and multiple threads are exploited. On the graphics
unit, the work is organised into tiles with shared memory. The neural
processing unit implementation was written entirely at the level of the
MLIR-AIE framework. The implementations were verified over 2700 cases in total.

\paragraph{Measurement harness.} Each case is measured in two explicitly
distinct scopes, algorithmic and end-to-end. Energy consumption is recorded
with a common methodology on all three architectures, the volume of data moved
is measured directly, and the performance of each run is analysed against the
roofline model derived for the same machine. This methodology is independent of
the top-k problem and can be applied to any workload considered for offloading
to an accelerator.

\paragraph{Findings.} Top-k selection is a data-movement problem on every
architecture examined, and no implementation is limited by computational
capability. The ranking of the two algorithms depends on the ratio $n/k$ and
stays the same across all architectures. The measurement scope completely
overturns the ranking of the architectures. Energy-wise, the neural processing
unit draws the lowest power of all implementations, and yet ends up with a
particularly high energy footprint, simply because it needs far more time.

\paragraph{Reproducibility.} The entire experimental procedure runs from
scripts included in the repository. The measurement conditions are recorded
alongside the results. Details are given in appendix~\ref{app:reproducibility}.

\section{Structure of the work}
\label{sec:intro-outline}

\emph{Chapter~\ref{ch:background}} defines the problem, presents the algorithm
families that solve it, describes the execution models of the three
architectures and concludes with the gap this work is called to fill.

The methodology follows (\emph{chapter~\ref{ch:methodology}}). It presents what
is measured, with what semantics, against which baselines and under what
conditions. It is a prerequisite for reading the results of
chapter~\ref{ch:evaluation}.

Next come the implementation chapters, one per architecture
(\emph{chapters~\ref{ch:cpu}, \ref{ch:gpu}, \ref{ch:npu}}), with a common
internal structure so that corresponding sections can be compared directly:
design choices, the two algorithms, memory management, the reference
implementation, limiting factors.

\emph{Chapter~\ref{ch:evaluation}} presents and interprets the results, namely
execution time, the comparison of architectures, memory footprint and energy
consumption, for different parameters, data types and input distributions.

\emph{Chapter~\ref{ch:conclusions}} closes with the conclusions, the limits
within which they hold and possible directions for continuation.

Acronyms and reproduction instructions are found in the appendices.
